\scanpage{53}
donc $\operatorname{Drap}(C)$ est un \uline{torseur à droite} sous $G_0=\operatorname{Aut}(\unsure{C^{0}|C_0})$ (et un torseur à gauche sous $G=\operatorname{Aut}(C)$, qui est le commutant de $G_0$\ldots)

Le foncteur
\diagram{\begin{tikzcd}[ampersand replacement=\&,row sep=1ex,column sep=large]
C\arrow[r] \& \operatorname{Drap}(C)\\
n\text{-Cubes}\arrow[r] \& G_0\text{-tors. à dr}
\end{tikzcd}}
est une équivalence de catégories

\uline{NB} $(C_0,\unsure{R_0|D_0})$ est défini à isom unique près -- donc aussi $G_0$. Si on veut, on prend pour $C_0$ le cube \unsure{standart|standard} $C_n$, \deleted{\unsure{avec le drapeau}} \inserted{qui du pt de} vue combinatoire correspond à $B=[1,n]\times\{\pm1\}$, avec comme drapeau celui défini par la section cte de valeur $1$, et ses restrictions aux sous-intervalles $[1,i]$ $(0\leq i\leq n)$. Ainsi on trouve un isom. can.
\[
\begin{aligned}
G_0&\simeq\{\pm1\}^n\cdot\unsure{\mathfrak S_n|\mathscr S_n}
\qquad\text{(produit }\tfrac12\text{ direct}\\
&\simeq\operatorname{Aut}(C_n)
\end{aligned}
\]
% archive sheet number: 52
